% Ad Familiares 4.2 --- headnote
% ad-familiares-04-02, 28 April 49 BC, Cumae

\textit{Cicero to Servius Sulpicius Rufus, written from his
villa at Cumae on the fourth day before the Kalends of May ---
Perseus dateline \textit{Scr. in Cumano a. d. iv K. Mai. a.
705 (49)}, that is, 28 April 49 BC. Caesar has crossed the
Rubicon in January, Pompey has evacuated Italy to Greece in
March, and Cicero --- having stayed behind --- is now agonising
on the Campanian coast over whether to follow Pompey, to
sit out the war, or to make some accommodation with the
victor. Servius Sulpicius, the great jurist and Cicero's
consular colleague in spirit if not in year, is on
neighbouring property and has been weighing the same
question; his wife Postumia and his son Servius (``our friend
Servius'' in the body of the letter) have come over to Cicero
in person to press for a face-to-face conversation.}

\textit{The note is short, but its three middle propositions
are among the sharpest things Cicero ever wrote about the
problem of conduct under tyranny: that what is upright is
plain and what is expedient is obscure; that the honest man
will admit nothing as expedient that is not also honourable;
and that, of the two remaining options for the citizen who
cannot endorse what is happening, mere acquiescence is shameful
and participation is dangerous. The conclusion --- ``one must
withdraw'' (\textit{discedendum}) --- is the closest Cicero has
come on paper to deciding for Pompey; the practical question
of \textit{where} to go is left open for the conversation
Servius is being invited to have with him at Cumae.}
